% Единые справочные данные для тренировочных вариантов.
\noindent Ниже приведены справочные данные, которые могут понадобиться при выполнении работы.
\begin{center}\bfseries Десятичные приставки\end{center}
\noindent\begin{tabularx}{\linewidth}{|X|c|c|X|c|c|}\hline
Наименование & Обозначение & Множитель & Наименование & Обозначение & Множитель\\\hline
гига & Г & $10^9$ & санти & с & $10^{-2}$\\\hline
мега & М & $10^6$ & милли & м & $10^{-3}$\\\hline
кило & к & $10^3$ & микро & мк & $10^{-6}$\\\hline
гекто & г & $10^2$ & нано & н & $10^{-9}$\\\hline
деци & д & $10^{-1}$ & пико & п & $10^{-12}$\\\hline
\end{tabularx}
\par\smallskip
\noindent\begin{tabularx}{\linewidth}{|X|>{\raggedright\arraybackslash}p{.42\linewidth}|}\hline
\multicolumn{2}{|l|}{\textbf{\itshape Константы}}\\\hline
число $\pi$ & $\pi=3{,}14$\\
ускорение свободного падения на Земле & $g=10\,\text{м/с}^2$\\
гравитационная постоянная & $G=6{,}67\cdot10^{-11}\,\text{Н}\cdot\text{м}^2/\text{кг}^2$\\
универсальная газовая постоянная & $R=8{,}31\,\text{Дж}/(\text{моль}\cdot\text{К})$\\
постоянная Больцмана & $k=1{,}38\cdot10^{-23}\,\text{Дж/К}$\\
постоянная Авогадро & $N_A=6\cdot10^{23}\,\text{моль}^{-1}$\\
скорость света в вакууме & $c=3\cdot10^8\,\text{м/с}$\\
коэффициент пропорциональности в законе Кулона & $k=\dfrac{1}{4\pi\varepsilon_0}=9\cdot10^9\,\text{Н}\cdot\text{м}^2/\text{Кл}^2$\\
элементарный электрический заряд & $e=1{,}6\cdot10^{-19}\,\text{Кл}$\\
постоянная Планка & $h=6{,}6\cdot10^{-34}\,\text{Дж}\cdot\text{с}$\\\hline
\end{tabularx}
\par\smallskip
\noindent\begin{tabularx}{\linewidth}{|X|X|}\hline
\multicolumn{2}{|l|}{\textbf{\itshape Соотношения между единицами}}\\\hline
температура & $0\,\text{К}=-273\,^{\circ}\text{C}$\\
атомная единица массы & $1\,\text{а.е.м.}=1{,}66\cdot10^{-27}\,\text{кг}$\\
энергия одной атомной единицы массы & $1\,\text{а.е.м.}\,c^2=931{,}5\,\text{МэВ}$\\
электронвольт & $1\,\text{эВ}=1{,}6\cdot10^{-19}\,\text{Дж}$\\\hline
\end{tabularx}
\par\smallskip
\noindent\begin{tabularx}{\linewidth}{|X|X|}\hline
\multicolumn{2}{|l|}{\textbf{\itshape Масса частиц}}\\\hline
электрона & $9{,}1\cdot10^{-31}\,\text{кг}\approx5{,}5\cdot10^{-4}\,\text{а.е.м.}$\\
протона & $1{,}673\cdot10^{-27}\,\text{кг}\approx1{,}007\,\text{а.е.м.}$\\
нейтрона & $1{,}675\cdot10^{-27}\,\text{кг}\approx1{,}008\,\text{а.е.м.}$\\\hline
\end{tabularx}
\par\smallskip
\noindent\begin{tabularx}{\linewidth}{|X|r|X|r|}\hline
\multicolumn{4}{|l|}{\textbf{\itshape Плотность}}\\\hline
воды & $1000\,\text{кг/м}^3$ & подсолнечного масла & $900\,\text{кг/м}^3$\\
древесины (сосна) & $400\,\text{кг/м}^3$ & алюминия & $2700\,\text{кг/м}^3$\\
керосина & $800\,\text{кг/м}^3$ & железа & $7800\,\text{кг/м}^3$\\
 & & ртути & $13600\,\text{кг/м}^3$\\\hline
\end{tabularx}
\noindent\begin{tabularx}{\linewidth}{|X|r|X|r|}\hline
\multicolumn{4}{|l|}{\textbf{\itshape Удельная теплоёмкость}}\\\hline
воды & $4200\,\text{Дж/(кг}\cdot\text{К)}$ & алюминия & $900\,\text{Дж/(кг}\cdot\text{К)}$\\
льда & $2100\,\text{Дж/(кг}\cdot\text{К)}$ & меди & $380\,\text{Дж/(кг}\cdot\text{К)}$\\
железа & $460\,\text{Дж/(кг}\cdot\text{К)}$ & чугуна & $500\,\text{Дж/(кг}\cdot\text{К)}$\\\hline
свинца & $130\,\text{Дж/(кг}\cdot\text{К)}$ & & \\\hline
\end{tabularx}
\clearpage
\par\smallskip
\noindent\begin{tabularx}{\linewidth}{|X|r|}\hline
\multicolumn{2}{|l|}{\textbf{\itshape Удельная теплота}}\\\hline
парообразования воды & $2{,}3\cdot10^6\,\text{Дж/кг}$\\
плавления свинца & $2{,}5\cdot10^4\,\text{Дж/кг}$\\
плавления льда & $3{,}3\cdot10^5\,\text{Дж/кг}$\\\hline
\end{tabularx}
\par\smallskip
\noindent\fbox{\begin{minipage}{.96\linewidth}\textbf{\itshape Нормальные условия:} давление $10^5\,\text{Па}$, температура $0\,^{\circ}\text{C}$.\end{minipage}}
\par\smallskip
\noindent\begin{tabularx}{\linewidth}{|X|X|X|X|}\hline
\multicolumn{4}{|l|}{\textbf{\itshape Молярная масса}}\\\hline
азота & $28\cdot10^{-3}$ & гелия & $4\cdot10^{-3}$\\
аргона & $40\cdot10^{-3}$ & кислорода & $32\cdot10^{-3}$\\
водорода & $2\cdot10^{-3}$ & лития & $6\cdot10^{-3}$\\
воздуха & $29\cdot10^{-3}$ & неона & $20\cdot10^{-3}$\\
воды & $18\cdot10^{-3}$ & углекислого газа & $44\cdot10^{-3}$\\\hline
\end{tabularx}
\noindent\small\hfill Единицы молярной массы: кг/моль.\normalsize
